% ---
% RESUMOS
% ---

% RESUMO em português
\setlength{\absparsep}{18pt} % ajusta o espaçamento dos parágrafos do resumo
\begin{resumo}
 O projeto tem como objetivo analisar o consumo de energia da Java Virtual Machine (JVM) através da comparação de diversos \textit{benchmarks}. Utilizando-se dos \textit{benchmarks} DaCapo e da ferramenta JRAPL, foram realizados uma série de testes com diferentes frequências do processador e configurações de tamanho de \textit{heap}. A partir das métricas de consumo de energia e tempo de execução, foi possível analisar o comportamento da energia para diferentes cargas de trabalho. Os resultados mostram que o perfil de consumo de energia é similar para todos os \textit{benchmarks}, ainda que haja diferenças no perfil de consumo de processador e memória. Ele demonstra que o consumo de energia é menor para tamanhos de heap maiores que 3 vezes o mínimo necessário e frequências de processador intermediárias.

 \textbf{Palavras-chaves}: Java. JVM. eficiência energética. consumo de energia. computação verde
\end{resumo}

% ABSTRACT in english
\begin{resumo}[Abstract]
 \begin{otherlanguage*}{english}
   This project aims to analyze the energy consumption of the Java Virtual Machine (JVM) by comparing various benchmarks. Using the DaCapo benchmarks and the JRAPL tool, a series of tests were conducted with different processor frequencies and heap size configurations. Based on energy consumption and execution time metrics, it was possible to analyze energy behavior for different workloads. The results show that the energy consumption profile is similar for all \textit{benchmarks}, even though there are differences in processor and memory usage profiles. It demonstrates that energy consumption is lower for heap sizes larger than 3 times the minimum required and at intermediate processor frequencies.

   \vspace{\onelineskip}
 
   \noindent 
   \textbf{Keywords}: Java. JVM. energy efficiency. energy consumption. green computing
 \end{otherlanguage*}
\end{resumo}